\documentclass[11pt,a4paper]{article}

\usepackage[margin=2.5cm]{geometry}
\usepackage{amsmath,amssymb,bm}
\usepackage{booktabs}
\usepackage{microtype}
\usepackage{hyperref}

\title{Comparison of Inverse, Forward, and Physics-Based Fitting Methods for Ion Beam Analysis}
\author{}
\date{}

\begin{document}

\maketitle

\begin{abstract}
This study compares three approaches for the analysis of ion beam analysis (IBA) spectra:
(1) a direct inverse model that predicts an elemental depth profile from a measured spectrum,
(2) fitting based on a learned forward model, and
(3) conventional fitting using SIMNRA.
The present document describes the probabilistic formulation of the inverse model.
The forward-model and SIMNRA-based fitting approaches are included as placeholders and will be described separately.
\end{abstract}

\section{Problem formulation}

Let
\[
\bm{S}
\]
denote an experimental or simulated IBA spectrum. The quantity to be inferred is the elemental areal-density map
\[
\bm{A}
\in
\mathbb{R}_{+}^{N_{\mathrm{lay}}\times N_{\mathrm{el}}},
\]
where $N_{\mathrm{lay}}$ is the maximum number of layers and
$N_{\mathrm{el}}$ is the maximum number of elements.

The entry
\[
A_{l,e}
\]
denotes the areal density of element $e$ in layer $l$. A value of zero indicates that the corresponding element is absent from that layer.

The goal of the inverse model is therefore to approximate the posterior distribution
\[
p(\bm{A}\mid\bm{S}).
\]

Rather than predicting only a single deterministic elemental depth profile, the model predicts both the probability that a given element is present and the conditional probability distribution of its areal density.

\section{Method 1: Probabilistic inverse model}

\subsection{Latent elemental-existence mask}

We introduce a binary elemental-existence matrix
\[
\bm{M}
\in
\{0,1\}^{N_{\mathrm{lay}}\times N_{\mathrm{el}}},
\]
with
\[
M_{l,e}
=
\begin{cases}
1, & \text{if element $e$ is present in layer $l$},\\
0, & \text{otherwise}.
\end{cases}
\]

The posterior distribution of the areal-density map can then be written by marginalizing over all possible elemental-existence configurations:
\[
\boxed{
p(\bm{A}\mid\bm{S})
=
\sum_{\bm{M}}
p(\bm{A}\mid\bm{M},\bm{S})\,
p(\bm{M}\mid\bm{S})
}
\label{eq:posterior_decomposition}
\]

This decomposition separates the inverse problem into two related tasks:
\begin{enumerate}
    \item determining which elements are present in each layer, and
    \item determining their areal densities conditional on their presence.
\end{enumerate}

\subsection{Element-existence probabilities}

The neural network predicts an elemental-existence probability matrix
\[
\bm{P}
\in
[0,1]^{N_{\mathrm{lay}}\times N_{\mathrm{el}}},
\]
where
\[
P_{l,e}
=
p(M_{l,e}=1\mid\bm{S}).
\]

Under the simplifying assumption that the individual existence variables are conditionally independent given the spectrum,
\[
\boxed{
p(\bm{M}\mid\bm{S})
=
\prod_{l=1}^{N_{\mathrm{lay}}}
\prod_{e=1}^{N_{\mathrm{el}}}
P_{l,e}^{M_{l,e}}
(1-P_{l,e})^{1-M_{l,e}}
}
\label{eq:mask_probability}
\]

Thus, the network does not need to predict probabilities for all
$2^{N_{\mathrm{lay}}N_{\mathrm{el}}}$ possible elemental configurations explicitly.
Instead, it predicts one Bernoulli probability for each layer--element pair.

\subsection{Conditional areal-density distribution}

Conditional on the existence of an element, its areal density is modeled by a Gaussian distribution.

The network predicts two additional matrices:
\[
\bm{Y}
\in
\mathbb{R}^{N_{\mathrm{lay}}\times N_{\mathrm{el}}},
\]
and
\[
\bm{R}
\in
\mathbb{R}_{+}^{N_{\mathrm{lay}}\times N_{\mathrm{el}}},
\]
where $Y_{l,e}$ is the predicted mean areal density and $R_{l,e}$ is the predicted standard deviation.

For a present element,
\[
A_{l,e}\mid M_{l,e}=1,\bm{S}
\sim
\mathcal{N}
\left(
Y_{l,e},
R_{l,e}^{2}
\right).
\]

For an absent element, the areal density is exactly zero. Therefore,
\[
p(A_{l,e}\mid M_{l,e},\bm{S})
=
\begin{cases}
\delta(A_{l,e}),
&
M_{l,e}=0,
\\[4pt]
\mathcal{N}
\left(
A_{l,e};
Y_{l,e},
R_{l,e}^{2}
\right),
&
M_{l,e}=1,
\end{cases}
\]
where $\delta(\cdot)$ denotes a Dirac point mass at zero.

Assuming conditional independence between layer--element entries,
\[
\boxed{
p(\bm{A}\mid\bm{M},\bm{S})
=
\prod_{l,e}
\left[
\mathcal{N}
\left(
A_{l,e};
Y_{l,e},
R_{l,e}^{2}
\right)
\right]^{M_{l,e}}
\left[
\delta(A_{l,e})
\right]^{1-M_{l,e}}
}
\label{eq:conditional_density}
\]

\subsection{Resulting predictive posterior}

Substituting Eqs.~\eqref{eq:mask_probability} and
\eqref{eq:conditional_density} into Eq.~\eqref{eq:posterior_decomposition}
and marginalizing the binary existence variables yields
\[
\boxed{
p(\bm{A}\mid\bm{S})
=
\prod_{l,e}
\left[
(1-P_{l,e})\,\delta(A_{l,e})
+
P_{l,e}\,
\mathcal{N}
\left(
A_{l,e};
Y_{l,e},
R_{l,e}^{2}
\right)
\right].
}
\label{eq:final_posterior}
\]

Hence, for every layer--element pair, the predictive distribution consists of
\begin{itemize}
    \item a point mass at zero with probability $1-P_{l,e}$, representing absence of the element, and
    \item one Gaussian component with probability $P_{l,e}$, representing the areal-density distribution conditional on presence.
\end{itemize}

This can be interpreted as a spike-and-slab form of a mixture-density model. The continuous part contains a single Gaussian mixture component, while the additional discrete component represents elemental absence.

\subsection{Neural-network outputs}

For an input spectrum $\bm{S}$, the inverse network therefore produces three matrices:
\[
\boxed{
\bm{Y},\qquad
\bm{R},\qquad
\bm{P}
}
\]
all of shape
\[
N_{\mathrm{lay}}\times N_{\mathrm{el}}.
\]

\begin{table}[h]
\centering
\begin{tabular}{lll}
\toprule
Output & Shape & Interpretation \\
\midrule
$\bm{Y}$ &
$N_{\mathrm{lay}}\times N_{\mathrm{el}}$ &
Mean areal density conditional on presence \\
$\bm{R}$ &
$N_{\mathrm{lay}}\times N_{\mathrm{el}}$ &
Predicted standard deviation \\
$\bm{P}$ &
$N_{\mathrm{lay}}\times N_{\mathrm{el}}$ &
Probability of elemental presence \\
\bottomrule
\end{tabular}
\caption{Outputs of the probabilistic inverse model.}
\label{tab:outputs}
\end{table}

In practice, the network may predict an unconstrained quantity
$r_{l,e}$ and transform it into a positive standard deviation using, for example,
\[
R_{l,e}
=
\operatorname{softplus}(r_{l,e})+\epsilon,
\]
where $\epsilon>0$ is a small numerical constant.

Similarly, elemental-existence probabilities can be obtained from logits $q_{l,e}$ using
\[
P_{l,e}
=
\operatorname{sigmoid}(q_{l,e}).
\]

\subsection{Training objective}

For supervised training, the true areal-density matrix
$\bm{A}^{\mathrm{true}}$ is known. The corresponding true existence mask is
\[
M_{l,e}^{\mathrm{true}}
=
\begin{cases}
1, & A_{l,e}^{\mathrm{true}}>0,\\
0, & A_{l,e}^{\mathrm{true}}=0.
\end{cases}
\]

The network is trained by minimizing the negative log-likelihood
\[
\mathcal{L}
=
-\log
p\left(
\bm{A}^{\mathrm{true}}
\mid
\bm{S}
\right).
\]

For an absent element,
\[
M_{l,e}^{\mathrm{true}}=0,
\]
the contribution to the loss is
\[
\boxed{
\mathcal{L}_{l,e}^{\mathrm{absent}}
=
-\log(1-P_{l,e}).
}
\]

For a present element,
\[
M_{l,e}^{\mathrm{true}}=1,
\]
the contribution is
\[
\mathcal{L}_{l,e}^{\mathrm{present}}
=
-\log P_{l,e}
-
\log
\mathcal{N}
\left(
A_{l,e}^{\mathrm{true}};
Y_{l,e},
R_{l,e}^{2}
\right).
\]

Expanding the Gaussian term gives
\[
\boxed{
\mathcal{L}_{l,e}^{\mathrm{present}}
=
-\log P_{l,e}
+
\log R_{l,e}
+
\frac{
\left(
A_{l,e}^{\mathrm{true}}-Y_{l,e}
\right)^2
}{
2R_{l,e}^{2}
}
+
\frac{1}{2}\log(2\pi).
}
\]

The complete loss can therefore be written as
\[
\boxed{
\begin{aligned}
\mathcal{L}
=
\sum_{l,e}
\Bigg[
&
-\left(1-M_{l,e}^{\mathrm{true}}\right)
\log\left(1-P_{l,e}\right)
-M_{l,e}^{\mathrm{true}}\log P_{l,e}
\\
&
+
M_{l,e}^{\mathrm{true}}
\left(
\log R_{l,e}
+
\frac{
\left(
A_{l,e}^{\mathrm{true}}-Y_{l,e}
\right)^2
}{
2R_{l,e}^{2}
}
+
\frac{1}{2}\log(2\pi)
\right)
\Bigg].
\end{aligned}
}
\label{eq:loss}
\]

Thus, elemental classification, areal-density regression, and uncertainty estimation are optimized jointly through a single probabilistically defined likelihood.

\subsection{Relation to mixture-density networks}

A conventional mixture-density network represents a conditional density as
\[
p(y\mid x)
=
\sum_{k=1}^{K}
\pi_k(x)
\,
\mathcal{N}
\left(
y;
\mu_k(x),
\sigma_k^2(x)
\right).
\]

In the present formulation, the number of continuous Gaussian components is $K=1$. Each layer--element pair additionally contains a discrete zero component describing elemental absence:
\[
p(A_{l,e}\mid\bm{S})
=
(1-P_{l,e})\delta(A_{l,e})
+
P_{l,e}
\mathcal{N}
\left(
A_{l,e};
Y_{l,e},
R_{l,e}^{2}
\right).
\]

The model can therefore be regarded as a single-component MDN augmented with an explicit sparse structural component.

\subsection{Use for direct inference and optimizer initialization}

The probabilistic output can be used directly to obtain an estimated elemental depth profile. It can also be used to generate multiple plausible initial solutions for subsequent numerical fitting.

In particular, samples can be drawn from
\[
p(\bm{A}\mid\bm{S})
\]
and passed as multiple initial guesses to an optimizer. This allows the inverse network to act either as a standalone predictor or as an informed initializer for physics-based refinement.

\section{Method 2: Learned forward model}

\textit{Placeholder.}

This section will describe fitting in which the computationally expensive physics-based forward calculation is replaced by a learned surrogate forward model. The surrogate predicts an IBA spectrum for a given elemental depth profile, while numerical optimization is used to solve the inverse problem.

\section{Method 3: SIMNRA-based fitting}

\textit{Placeholder.}

This section will describe conventional fitting in which SIMNRA is used directly as the forward model and the elemental depth profile is determined by numerical optimization.

\section{Comparison strategy}

The three approaches will ultimately be compared in terms of quantities such as:
\begin{itemize}
    \item accuracy of the recovered elemental depth profile,
    \item agreement between measured and reconstructed spectra,
    \item computational cost,
    \item number of forward evaluations required,
    \item robustness as the number of layers and elements increases,
    \item sensitivity to weak or ambiguous elemental signals, and
    \item uncertainty or repeatability of the inferred solution.
\end{itemize}

The precise evaluation protocol and benchmark metrics will be defined separately.

\end{document}
